% Game sections: 5..11 (the credits are 12, see override-en.tex).

\sect{5. How to play}

\textbf{The story.} Something evil hides beneath the city of Waterdeep. Piergeiron, the first Lord of the city, gathers a party of adventurers and gives them a commission: go down into the sewers, find out what threatens the city and destroy it. As they enter, a cave-in closes the way back. From there the only way is down: the sewers, the dwarven ruins, the caves of the dark elves and, at the bottom, the lair of the beholder Xanathar.

\textbf{The goal.} Get through the 12 levels of the dungeon, stay alive and put an end to Xanathar.

\textbf{Everything is seen through the party's eyes.} The large window shows what lies ahead. You move square by square and turn in steps of 90 degrees. Walls pass by, and you see items on the floor, doors, levers, buttons, writings and monsters coming close to fight. Monsters do not wait: if you stand still, they keep moving and attacking.

\textbf{The party.} You start with four characters. The two at the top are the \textbf{front rank}: they are the only ones who can fight in melee (swords, maces, axes). The two below are the \textbf{rear rank}: they can only attack with ranged weapons (bows, thrown daggers or darts) or with spells, but remain exposed, especially to attacks from the sides or behind. On the way you will meet other adventurers who may join the party, up to six characters. To \textbf{change the order}, right-click on a character's name (\texttt{Swapping} shows in its place), then right-click on another name: the two change places. A click anywhere else on the panel cancels it.

\textbf{Take, use, select.} The key to the game is the difference between these three actions:
\begin{itemize}
\item \textbf{Take} (left button): pick up or put down items. What you take becomes the mouse pointer, and you put it wherever you want.
\item \textbf{Use} (right button on a hand, or the keys \key{1}--\key{6} and \key{Q}--\key{Y}): attack with a weapon, cast a spell with the spell book or the holy symbol, drink a potion, read a scroll\dots
\item \textbf{Select} (left button): press menu buttons, choose a spell from a list, point at a character.
\end{itemize}

\textbf{Examining the world.} Left-click on a writing, a drain or a body: what the party sees appears in the message area. The same click works levers, buttons and chains.

\textbf{Doors.} Not all of them open the same way. Many have a button or a lever beside them (click on it, or \key{SPACE} with the door ahead). Locked doors open by taking the right key and clicking with it on the lock; a thief can try the same with his lockpicks. A half-open door can be forced by clicking on its lower part: the strongest character makes the attempt.

\textbf{Stairs and portals} lead to other levels. The game remembers what you left on each level (open doors, dead monsters, dropped items) in case you come back.

\sect{6. Making the party}

A new game starts on the \textbf{Character Generation} screen, where you make the four adventurers. Everything is done with the mouse, as in the original; on the keyboard the cursor keys choose a box and \key{RETURN} opens it.

\begin{center}
\includegraphics[width=0.48\textwidth]{\shotdir chargen-race.png}\hfill
\includegraphics[width=0.48\textwidth]{\shotdir chargen-stats.png}
\end{center}

\begin{enumerate}
\item \textbf{Choose an empty box} of the four on the left.
\item \textbf{Race and sex, class and alignment}, in that order. Only the classes each race allows are listed; \texttt{BACK} returns to the previous step. A party with a paladin takes no evil characters, and the other way round.
\item \textbf{The abilities} are rolled: strength, intelligence, wisdom, dexterity, constitution and charisma, with the armour class (\texttt{AC}), hit points (\texttt{HP}) and level. Before that, choose a \textbf{portrait}: four are shown, and the arrows show others.
\item With the buttons on the right: \texttt{REROLL} rolls again; \texttt{MODIFY} lets you raise or lower each value with \texttt{+} and \texttt{-} within the limits of the race and class (the one being changed is red; the up and down cursor keys go to another; \texttt{OK} ends); \texttt{FACES} changes the portrait; \texttt{KEEP} keeps the character.
\item \textbf{Type the name} (up to ten letters) and press \key{RETURN}.
\end{enumerate}

Click a box with a character to see its abilities: \texttt{DELETE} removes it so that you can make another one there. With all four made, \texttt{PLAY} appears: click it (or press \key{P}) and each character receives the equipment of its class (armour, weapons, the mage's spell book, the cleric's holy symbol, food). The first one also carries the commission letter of Waterdeep.

\subsection*{Choosing a balanced party}
For a first game, combine two front-line fighters with a cleric and a mage behind them. The cleric brings healing and food; the mage brings attacks and protections. A thief or a multi-class character adds lockpicks without necessarily giving up another role. Carry ranged weapons for the rear rank.

\begin{reference}{Class}{Role in the party}
Fighter & Weapons and armour specialist; a good choice for taking attacks from the front.\\
Ranger & A combatant sharing many equipment options with the fighter. Do not assume a ranger has a mage's spell book.\\
Paladin & A good-aligned combatant with Lay on Hands and access to a selection of prayers. Alignment restricts the party's composition.\\
Cleric & Healing, food, protections and help against undead. Can wear armour, but has a more restricted choice of weapons.\\
Mage & Learns spells from a book and scrolls. Physical weapons and protection are limited: keep the mage in the rear rank.\\
Thief & Uses lockpicks on locks. Not every lock can be opened this way; keep keys as well.\\
Multiple classes & Combines the roles of the listed classes. Experience is shared between them, so progress differs from a single-class character.\\
\end{reference}

\textbf{Races.} Humans, elves, half-elves, dwarves, gnomes and halflings offer different class combinations and ability limits. Character generation shows the actual choices available for each selection. An elf cannot receive this version's Raise Dead spell; consider this before making the party. Sex and portrait define the character without changing the controls.

\begin{reference}{Ability}{What to look for}
Strength (STR) & Important for physical combat and forcing doors. Some fighters also display exceptional strength alongside 18.\\
Intelligence (INT) & The mage's ability; consider it when choosing a character who will study magic.\\
Wisdom (WIS) & The cleric's ability; relevant to preparing prayers.\\
Dexterity (DEX) & Helps protect the character. Check its effect on AC, especially with little armour.\\
Constitution (CON) & Important for endurance and hit points. Do not neglect the front rank's HP.\\
Charisma (CHA) & Part of the character sheet and some class requirements. It is not a negotiation command and does not replace dialogue choices.\\
\end{reference}

\textbf{Alignment.} Describes the character's moral outlook and restricts some classes and party combinations. It does not add an ability button. Review the characters before pressing PLAY; you need not accept the first roll.


\sect{7. The adventure screen}

\begin{center}
\includegraphics[width=0.86\textwidth]{\imgdir hud.png}
\end{center}

\begin{tabularx}{\textwidth}{@{}clX@{}}
\toprule
\textbf{\tNo} & \textbf{\tElement} & \textbf{\tTellsYou}\\
\midrule
1 & View & What lies ahead of the party: walls, doors, items and monsters. Left click: take, throw, put down or examine.\\
2 & Portrait and name & Who each character is. Click: the inventory. Greyed: unconscious; a skull: dead. When hit, a red splash shows the damage.\\
3 & Upper hand & What the character holds in the primary hand, usually the weapon. Right click or \key{1}--\key{6}: use it. After use it stays grey for a moment.\\
4 & Hit point bar & The hit points: green; yellow when less than a third is left; red when they reach zero and the character falls.\\
5 & Lower hand & The secondary hand: shield, second weapon, spell book, holy symbol\dots{} Right click or \key{Q}--\key{Y}: use it.\\
6 & Arrows & Step forward and back, turn, step sideways.\\
7 & Compass & Where the party is facing. While a spell is being cast, the spell menu appears here.\\
8 & Messages & What happens: hits, poison, writings of the dungeon, what other characters say.\\
9 & \texttt{CAMP} & The camp: rest, prepare spells, preferences, and save or load the game.\\
\bottomrule
\end{tabularx}

\sect{8. The inventory and the character sheet}

Click on a portrait and the right panel shows that character's inventory. Return to the game by clicking on the portrait again.

\begin{center}
\includegraphics[width=0.42\textwidth]{\imgdir inv.png}\hfill
\includegraphics[width=0.42\textwidth]{\shotdir stats-panel.png}
\end{center}

\begin{tabularx}{\textwidth}{@{}clX@{}}
\toprule
\textbf{\tNo} & \textbf{\tElement} & \textbf{\tTellsYou}\\
\midrule
1 & Portrait & Click: back to the adventure screen.\\
2 & \texttt{HP} and \texttt{FOOD} & Hit points and food. Without food, the character slowly loses hit points and cannot prepare spells.\\
3 & Plate & To eat: take a food ration and put it on the plate.\\
4 & Arrows & The previous or the next character (also useful to pass an item from one character to another).\\
5 & Backpack & Fourteen boxes for carrying things.\\
6 & Quiver & The arrows. Put arrows on it to fill it; the number tells how many are left.\\
7 & Primary hand & The weapon (the same upper hand as on the panel).\\
8 & Secondary hand & Shield, book, holy symbol, ammunition\dots\\
9 & Page & Turns to the character sheet (on the right): class, alignment, race, ability scores, experience and level.\\
10 & Body & Boxes for what is worn: head, neck, armour, wrists, rings and boots; on the right, the belt with three boxes for items close at hand (throwing weapons in it go to the hand by themselves when you throw).\\
\bottomrule
\end{tabularx}

To wear something, take it and put it in its box. If you carry an item on the pointer and click on a character's portrait, his inventory opens so you can store it there.

\subsection*{Reading the sheet and recognising conditions}
\textbf{AC} is armour class: lower is better, even below zero. \textbf{HP} shows current and maximum hit points. Experience leads to level gains; it is neither money nor a score to spend. Multi-class characters show their levels separately. A character's level and a spell's level are different things.

\begin{reference}{Condition}{What to do}
Wounded & Heal with magic, potions or safe rest. Food does not directly restore HP.\\
Hungry & Feed the character on the inventory plate. Hunger prevents spell recovery and eventually costs hit points.\\
Poisoned & Use an antidote or Neutralize Poison. Slow Poison only buys time; sleeping while poisoned can kill. Heal the damage already suffered separately.\\
Paralysed & The character cannot act; some attacks cause about two and a half minutes of paralysis. Protect them and use Remove Paralysis when available. Healing HP does not remove this condition.\\
Petrified & Cannot act. Healing does not reverse petrification. Keep an earlier save so you can recover the party if needed.\\
Unconscious & At zero HP or below, but above $-10$, the character has fallen. Healing above zero can return them to action.\\
Dead & At $-10$ HP the skull appears. Ordinary healing cannot help; Raise Dead, when applicable, returns the character with 1 HP. If the entire party falls, the game attempts to load the saved game.\\
\end{reference}

\textbf{The inventory does not pause the adventure.} Arrange equipment somewhere clear. Open camp to stop and think; resting advances time, so hunger, enemies and poison matter again.


\sect{9. Combat and magic}

\textbf{Attacking.} Use (right click, or \key{1}--\key{6}) the weapon in a front-rank character's hand when the monster is ahead, in the view. Ranged weapons are used the same way from either rank: a bow needs arrows in the other hand, and reloads from the quiver by itself as it shoots; daggers and darts are replaced from the belt. After each attack the hand stays grey until the character is ready again. When a monster dies, the party gains experience and the characters go up in level.

\textbf{Casting spells.} Mages use their \textbf{spell book}; clerics and paladins, their \textbf{holy symbol}. Use the book or the symbol (right click) and the spell menu appears in place of the compass:

\begin{center}
\includegraphics[width=0.6\textwidth]{\shotdir spells.png}
\end{center}

At the top are the spell levels (1 to 5); below, the prepared spells of that level. Click on one to cast it; if it affects a character, then click on that character's portrait. Attack spells can only be cast at monsters in sight. \texttt{ABORT SPELL} closes the menu. A cast spell is used up: it has to be prepared again by resting in the camp.

\textbf{Scrolls, potions and wands.} They are used like a weapon: in the hand and a right click. A scroll is used up when read. A mage can copy the spell of a scroll into his book (camp, \texttt{Scribe Scrolls}).

\textbf{Turning undead.} Clerics, and paladins from level 3 on, try by themselves to drive away the undead in front of them.

\textbf{Poison and paralysis.} Some monsters poison or paralyse. Poison takes hit points away little by little until it is cured (with a spell or a potion); the inventory page shows \texttt{POISONED}.

\subsection*{From the book to a ready spell: a complete example}
\begin{enumerate}
\item Enter CAMP and choose Memorize Spells for a mage, or Pray for Spells for a cleric or paladin.
\item Choose the character and spell level. Add copies with a left click or the + button; remove them with a right click or the $-$ button. If 0 remain, free a slot first. For example, replace one Aid with a Hold Person.
\item Leave the list and choose Rest Party. Selecting spells does not recover them immediately. Wait for messages confirming preparation; an interrupted rest may leave spells unfinished.
\item Return to the adventure, use the book or symbol from a hand, choose the level and a ready spell. If a recipient is requested, select that portrait.
\item Rest again to recover the spent spell. To learn new mage magic, first scribe its scroll and then select it for preparation.
\end{enumerate}

\textbf{Preparation keyboard controls.} Up and down highlight a spell; left and right change spell level. Return or Space adds one copy, the minus key removes one, and Esc exits keeping the selection. The number on the right is the selected quantity, not a field for typing digits.

\textbf{Touch, projectile and protection.} A projectile needs a clear direction towards its target; a touch spell may require the front rank or a free hand. With Shocking Grasp, Vampiric Touch and Flame Blade you must then use the hand containing the magic. A protection does not cause damage by itself. Section 13 explains each case.

\textbf{Combat rhythm.} Attack, wait for the hand to become available and use it again. Retreating or stepping sideways can buy time, provided you have checked the floor. Do not open the inventory expecting enemies to stop. Recover projectiles left on the floor once the area is safe.


\sect{10. The camp}

Click on \texttt{CAMP}. The game stops while you are in the camp.

\begin{center}
\includegraphics[width=0.32\textwidth]{\shotdir camp.png}\hfill
\includegraphics[width=0.32\textwidth]{\shotdir memorize.png}\hfill
\includegraphics[width=0.32\textwidth]{\shotdir options.png}
\end{center}

\begin{itemize}
\item \texttt{Rest Party}: the party rests hour by hour (the count shows in the window). After about eight hours each mage and cleric starts to regain their spells, one after another, and the message window tells about each. If somebody is injured and a cleric has healing spells, the game asks whether the cleric should heal the party. When no spell is left to regain but somebody is still injured, it asks whether to go on resting until healed. Every eight hours the characters regain a hit point and eat. A click interrupts the rest. You cannot rest with monsters near, and if one comes close the camp ends with a warning. Beware of poison: a poisoned character loses hit points every hour and may die while resting (the game warns first). If the party is hungry, the game asks whether to go on resting.
\item \texttt{Memorize Spells}: pick a mage and choose the spells to learn while resting. The numbers on the right tell how many of each there will be; \texttt{Clear} removes the choice and the buttons \texttt{1}--\texttt{5} change the level.
\item \texttt{Pray for Spells}: the same for clerics and paladins.
\item \texttt{Scribe Scrolls}: a mage copies into the spellbook the spell of a scroll in the inventory.
\item \texttt{Preferences}: \texttt{Sounds} turns the sound effects on and off. Music plays during the introduction.
\item \texttt{Game Options}: \texttt{Load Game}, \texttt{Save Game}, \texttt{Drop Character} (removes a member when there are five or six; details below) and \texttt{Quit Game} (back to MSX-DOS).
\item \texttt{Exit}: back to the game.
\end{itemize}

The keyboard works in the camp menus too: the arrows move the lit button, \texttt{RETURN} or the space bar choose it and \texttt{ESC} leaves (in the Yes/No questions the left and right arrows change the answer). When a character must be chosen, the arrows move the frame of changing colours and \texttt{RETURN} picks it. In the list of spells, the up and down arrows (or the pointer) light one, \texttt{RETURN}, the space bar or a click add it, the \texttt{-} key takes one away and the left and right arrows change the level.

\textbf{Saving.} There is a single saved game for each language: \texttt{Save Game} replaces the previous one. It is the one loaded by \texttt{LOAD GAME IN PROGRESS} on the title screen and when the whole party falls.

\subsection*{Every option and its buttons}
\begin{reference}{Menu / option}{Result}
Memorize / Pray & First select an eligible character. Buttons 1--5 change spell level, not character level. Left-click an entry to add a copy; right-click to remove one. The + and $-$ buttons act on the highlighted entry. With 0 slots remaining, remove a copy before adding another.\\
Clear & In spell preparation, clears the selection on the current spell-level page. It does not delete characters or remove learned spells from the book.\\
Exit preparation & Keeps the selection and returns to camp. You still need rest to recover pending copies.\\
Scribe Scrolls & Select a mage and a new spell they can copy. Mage scrolls carried by party members can be used. Scribing consumes the scroll and teaches the spell; it does not automatically prepare it. Divine scrolls cannot be copied.\\
Preferences: Sounds & Turns sound effects on or off when MoonSound is connected.\\
Preferences: Exit & Returns to camp.\\
Options: Load Game & Restores the current language's saved game. Progress since that save is lost. If loading fails, read the message before continuing.\\
Options: Save Game & Writes the current state and replaces the previous save for that language. Wait for completion; do not switch off while writing.\\
Options: Drop Character & Available only with five or six members. Select who leaves: it cannot be the first member, but is not restricted to recruited companions. Their possessions stay on the floor; review the party and recover what you need.\\
Options: Quit Game & Asks for confirmation before returning to MSX-DOS. Yes quits; No cancels. It does not save automatically.\\
Options: Exit & Returns to camp. Exit in CAMP's main menu resumes the adventure.\\
\end{reference}

\textbf{Questions during rest.} Accepting healing permits available cures to be used; continuing rest to heal takes longer and consumes more food. Accepting rest while hungry or poisoned does not remove those dangers. A click interrupts the process; monsters can also prevent or interrupt it.

\textbf{Dialogues.} Read the message and choose an offered answer: accepting or refusing help, recruiting or answering a character may change what happens. Not every encounter is a fight. A party of six needs a vacancy to recruit someone else.

\textbf{What is saved.} Each language has one save for the party and its progress. Choosing \texttt{Save Game} replaces that language's previous save. Starting a new party leaves the previous save available until you save again.

\sect{11. Equipment and adventure items}

The pictures show one example from each family; variants may share an icon. The colour of a potion or ring is not enough to identify its effect.

Each item has a place of use: hand, body, inventory or dungeon mechanism. The name alone does not establish its properties. Picking up a scroll with the left mouse button shows its spell name without consuming it. Unidentified potions and wands retain generic names: use Read Magic to identify them.

\textbf{Bonuses and curses.} A positive bonus can improve attack, damage or armour. A curse can worsen equipment and prevent removal. Save before testing suspicious items. There is no camp option for removing curses. Creatures which rust or destroy equipment can cost you items: carry replacements.

\refspace
\subsection*{Weapons and ammunition}

Base damage against small/medium and large enemies, respectively, before strength, bonuses and defences. \textbf{F}: fighter, ranger and paladin; \textbf{M}: mage; \textbf{C}: cleric; \textbf{T}: thief. Multiclass combinations can broaden the options; the game checks whether the character can use the equipment.

\begin{reference}{Name}{Effect, use and limits}
\itemname{axe}{Axe} & 1d8 / 1d10. One hand; F, T. Includes Drow Cleaver.\\[4pt]
\itemname{long-sword}{Long Sword} & 1d8 / 1d12. One hand; F, T. Variants include Slasher, Night Stalker, Severious and the adamantite sword.\\[4pt]
\itemname{short-sword}{Short Sword} & 1d6 / 1d8. One hand; F, T. Slicer belongs to this family.\\[4pt]
\itemname{dagger}{Dagger} & 1d4 / 1d3. F, M, T. Thrown when used from a hand; retrieve it afterwards. Backstabber, Flicka and Guinsoo are variants.\\[4pt]
\itemname{spear}{Spear} & 1d6 / 1d8. F, T. Used as a throwing weapon from a hand. It is not replaced unless you carry another.\\[4pt]
\itemname{halberd}{Halberd} & 1d10 / 2d6. Two hands; F, T. The Chieftain Halberd is a variant of this family.\\[4pt]
\itemname{mace}{Mace / Scepter} & 1d6 + 1 / 1d6. One hand; F, C, T. The Scepter of Kingly Might functions here as a weapon of this family; do not assume a special pushing power.\\[4pt]
\itemname{flail}{Flail} & 1d6 + 1 / 2d4. One hand; F, C, T. Used to strike from the front rank.\\[4pt]
\itemname{staff}{Staff} & 1d6 / 1d6. Two hands; all classes. A mage can wield it too.\\[4pt]
\itemname{bow}{Bow} & 1d6 / 1d6. Two hands; F, T. Needs arrows in the other hand and can reload from the quiver. Includes the drow bow.\\[4pt]
\itemname{sling}{Sling} & 1d4 / 1d4. Two hands; F, C, T. Uses rocks as ammunition. Check both hands if the no-ammunition warning appears.\\[4pt]
\itemname{dart}{Dart} & 1d3 / 1d2. F, M, T. A throwing weapon; can reload from the belt. Special darts have their own modifiers.\\[4pt]
\itemname{arrow}{Arrow and Rock} & Ammunition. Arrows go in the quiver and supply the bow. Rocks supply the sling and can also be thrown. Retrieve these items when they remain on the floor.\\[4pt]
\end{reference}

\refspace
\subsection*{Armour and protection}

Lower AC is better. Base values below exclude dexterity and magical items. Two differently named items can produce the same final AC.

\begin{reference}{Name}{Effect, use and limits}
\itemname{leather}{Leather Armour} & Base AC 8; F, C, T. An early choice for a character unable to wear heavy armour.\\[4pt]
\itemname{scale}{Scale Mail} & Base AC 6; F, C. Magical bonuses and dexterity may improve the result.\\[4pt]
\itemname{chain}{Chain Mail} & Base AC 5; F, C. Equip it on the body rather than in a hand.\\[4pt]
\itemname{banded}{Banded Armour} & Base AC 4; F, C. Compare the final result on the character sheet.\\[4pt]
\itemname{plate}{Plate Mail} & Base AC 3; F, C. Includes magical or cursed variants such as Plate Mail of Great Beauty; the name does not guarantee an improvement.\\[4pt]
\itemname{robe}{Robe} & All classes. A plain robe does not itself improve AC; the Robe of Defense has its own bonus.\\[4pt]
\itemname{shield}{Shield} & Reduces AC by 1 plus any magical bonus. Must occupy a usable hand; it does not combine with using a two-handed weapon.\\[4pt]
\itemname{bracers}{Bracers} & Their bonus may improve AC when worn on the wrists. Compare the final value after changing equipment.\\[4pt]
\itemname{ring-protection}{Ring of Protection} & Wear it on a finger. The better of two rings is used, and it does not add to magically enhanced armour.\\[4pt]
\itemname{helmet}{Helmet and Boots} & Equip them in the head and feet slots. Ordinary helmets and boots do not improve AC.\\[4pt]
\end{reference}

\refspace
\subsection*{Tools, keys and special items}

\begin{reference}{Name}{Effect, use and limits}
\itemname{spellbook}{Spell Book} & A mage needs it to prepare magic. Put it in a hand and right-click to open prepared spells. Casting does not consume the book.\\[4pt]
\itemname{holy-symbol}{Holy Symbol} & A cleric or paladin's tool. In a hand it opens prayers. Carrying it does not replace selecting and recovering spells through rest.\\[4pt]
\itemname{lockpicks}{Lockpicks} & Pick them up with the pointer and use them on a suitable lock. A thief can attempt to open it; failure does not necessarily mean the object is not a lock.\\[4pt]
\itemname{key}{Keys} & There are gold, silver, ruby, jeweled, skull, dwarven and drow keys. Bring the key to the matching slot or lock. A similar colour or shape does not guarantee a match.\\[4pt]
\itemname{gem}{Gems} & Keep gems: they may serve as offerings or mechanism pieces. Their use depends on the location; they are not currency to spend in a shop.\\[4pt]
\itemname{stone-object}{Stone Objects} & Stone dagger, holy symbol, medallion, necklace, orb, ring and scepter. These are travel-mechanism pieces. Their shape matters; keep each one and note where it was found.\\[4pt]
\itemname{bones}{Bones} & Human, dwarf and halfling remains can be connected to characters in the adventure. Keep them and pay attention to dialogue; they are not food.\\[4pt]
\itemname{letter}{Commission Letter and Notes} & Put a document in a hand and use it to read it. Notes provide context or clues. They are documents, not an automap that updates as you walk.\\[4pt]
\itemname{orb}{Orb of Power} & A special adventure item. Do not treat it as a rechargeable wand or ordinary weapon; keep unique objects until their purpose becomes clear.\\[4pt]
\itemname{necklace}{Necklaces and Medallions} & Worn on the neck. Some are ornaments, including the Luck Stone Medallion, which gives no combat bonus.\\[4pt]
\end{reference}

\sect{12. Consumables, wands and rings}

\textbf{Read a scroll name without spending it.} Pick it up with the left mouse button, from the floor or an inventory slot: the message area shows the scroll and its spell name. This works for mage and divine scrolls, even if the mage already knows the spell. Put it back in a slot with another left click. Right-clicking the hand uses and consumes it. Look up the spell name in this manual for its effect.

\textbf{Who receives the effect.} A potion is drunk by its holder. A scroll or wand which acts on a character applies to its holder when used; pass it to the intended recipient first. Prepared spells which request a target let you choose a portrait. Do not confuse the two procedures.

\textbf{Use has a cost.} Food, potions and scrolls disappear when consumed. Wands lose charges and disappear when exhausted, except Silvias, which keeps its charges. Rest does not recharge them. An attempt requiring a free hand may be refused until you free one.

After using a scroll or wand, wait for the hand to become available again.

\textbf{Identification and duration.} Pick up the item with the left button and cast Read Magic to identify it. Identified potions, wands and rings display their effect with their name. Clouds and timed effects survive saving and loading; rest advances their duration. Some effects end during rest, while longer effects may continue.

\begin{reference}{Name}{Effect, use and limits}
\itemname{rations}{Rations / Iron Rations} & Pick up food and drop it on the character inventory's plate. It is consumed and increases FOOD up to 100. Supplies come in different sizes and do not all feed equally.\\[4pt]
\itemname{egg}{Kenku Egg} & Also eaten on the plate. Its food contribution varies between eggs. Keep supplies for long rests.\\[4pt]
\itemname{potion-strength}{Potion of Giant Strength} & Sets Strength to 22 for about 150--240 seconds, then restores the previous value even after repeated doses. Does not heal HP or cure poison.\\[4pt]
\itemname{potion-healing}{Potion of Healing} & Restores 2d4 + 2 HP (4--10), up to the normal maximum. The character holding it drinks it.\\[4pt]
\itemname{potion-extra}{Potion of Extra Healing} & Restores 3d8 + 3 HP (6--27), up to the normal maximum. Does not revive the dead.\\[4pt]
\itemname{potion-poison}{Potion of Poison} & Can poison the drinker if the effect is not resisted. An unidentified bottle is not necessarily a cure: save before experimenting.\\[4pt]
\itemname{potion-vitality}{Potion of Vitality} & Fills FOOD to 100. Satisfies hunger but is not an HP healing potion.\\[4pt]
\itemname{potion-speed}{Potion of Speed} & Lets the drinker use their weapons twice as often for a few minutes.\\[4pt]
\itemname{potion-invisible}{Potion of Invisibility} & Makes the drinker invisible and harder for monsters to hit. It has no fixed expiry; using a hand or casting offensive magic can break it, as with the spell.\\[4pt]
\itemname{potion-antidote}{Cure Poison Potion} & Removes poisoning. Heal damage already suffered separately. The Dwarven Healing Potion also cures poison.\\[4pt]
\itemname{scroll-mage}{Mage Scroll} & Use from a hand to spend its spell once, or copy it into a spellbook through Scribe Scrolls. Scribing consumes the scroll too.\\[4pt]
\itemname{scroll-divine}{Divine Scroll} & Contains a single-use prayer. Cannot be copied into a mage's spellbook. When an item spell affects a character, using it applies the effect to its holder.\\[4pt]
\end{reference}

\refspace
\subsection*{Wands: available effects}

\begin{reference}{Name}{Effect, use and limits}
\itemname{wand}{Nonmagical Wand} & Produces no magical effect; the corresponding message appears.\\[4pt]
\itemname{wand}{Lightning} & Casts Lightning Bolt. Each valid use consumes a charge.\\[4pt]
\itemname{wand}{Frost} & Casts Cone of Cold ahead. Face the party correctly before using it.\\[4pt]
\itemname{wand}{Healing} & Uses Cure Serious Wounds (2d8 + 1 HP) on the holder. Pass it to the wounded character before activation.\\[4pt]
\itemname{wand}{Fireballs} & Casts Fireball. Needs a path towards the enemy.\\[4pt]
\itemname{wand}{Silvias / Slivias} & Pushes a monster in front of the party one square along an open passage into an empty square. Keeps its charges.\\[4pt]
\itemname{wand}{Magic Missiles} & Casts Magic Missile. A reserve of attacks when prepared spells run out.\\[4pt]
\end{reference}

\refspace
\subsection*{Rings: wear them to use them}

\begin{reference}{Name}{Effect, use and limits}
\itemname{ring-protection}{Protection} & Its bonus improves AC under the equipment rules in section 11. Keeping it in the backpack is not enough.\\[4pt]
\itemname{ring}{Sustenance} & Worn on a finger, prevents normal food loss, including during rest. Does not heal wounds.\\[4pt]
\itemname{ring}{Adornment} & An ornament. Gives no combat bonus.\\[4pt]
\itemname{ring}{Wizardry} & When worn on a finger, doubles available mage preparation slots for spell levels 4 and 5. Two rings do not multiply the bonus again. It grants neither access to higher spell levels nor knowledge of new spells.\\[4pt]
\itemname{ring}{Feather Fall} & When worn on a finger, prevents the falling damage marked for this protection. It does not protect against every trap or against ordinary attacks.\\[4pt]
\end{reference}

\sect{13. Spell and prayer reference}
The book marks mage spells and the holy symbol marks divine prayers. These are the pictures of their tools in the game, rather than a unique icon for each spell.

The tables distinguish \textbf{spell level} from character level. Preparing a spell requires access to that level, an available slot and, for a mage, knowledge of the spell in the book. Paladins have their own selection of prayers; the menu shows what they can prepare.

\textbf{Reading dice.} 1d8 means 1 to 8; 2d8 + 1 means rolling two eight-sided dice and adding one (3--17). These are base values: defences, resistances and saving throws can change the result. Healing never exceeds maximum HP. Do not assume repeated protection spells stack.

Multiclass characters use the level of the class casting the spell.

\refspace
\subsection*{Mage spells}

\refspace
\subsubsection*{Level 1}

\begin{reference}{Name}{Effect, use and limits}
\itemname{spellbook}{Armor} & One character. Provides base AC 6, adjusted for dexterity, when this improves protection. Ends after the recipient has taken damage totalling 8 + caster level; it does not absorb that damage.\\[4pt]
\itemname{spellbook}{Burning Hands} & Close attack against monsters in the square immediately ahead: 1d3 + 2 per caster level. Face the enemy before casting.\\[4pt]
\itemname{spellbook}{Detect Magic} & While active, magical item icons turn blue in inventory, hands and the cursor. Picking them up also adds [MAGIC] to the name. It does not identify them; use Read Magic to reveal their properties.\\[4pt]
\itemname{spellbook}{Magic Missile} & Forward projectile. Damage starts at 1d4 + 1 and scales with level through additional hits. Useful from the rear rank; walls stop its flight.\\[4pt]
\itemname{spellbook}{Read Magic} & Identifies the item held by the mouse cursor without consuming it, then displays its full name. Also works with potions, wands and rings. Available for any mage to prepare.\\[4pt]
\itemname{spellbook}{Shield} & Protects the caster. Caps AC at 4 when it was worse and stops incoming magic missiles. It does not grant immunity to all attacks or magic.\\[4pt]
\itemname{spellbook}{Shocking Grasp} & Requires a free hand. Creates a magical hand attack dealing 1d8 + level. Use that hand against a nearby enemy afterwards; selecting the spell does not launch a projectile.\\[4pt]
\end{reference}

\refspace
\subsubsection*{Level 2}

\begin{reference}{Name}{Effect, use and limits}
\itemname{spellbook}{Invisibility} & One character. Makes the character harder for monsters to hit, without granting invulnerability. Using a hand or casting offensive magic can break invisibility.\\[4pt]
\itemname{spellbook}{Knock} & Opens an ordinary door ahead and releases doors that could be physically forced. Does not alter walls, magical locks or puzzle mechanisms. Available to prepare for any mage with level-2 spell slots.\\[4pt]
\itemname{spellbook}{Melf's Acid Arrow} & Forward projectile: 2d4 for each complete three caster levels, with a minimum of one roll. Damage is applied on impact, without continuing damage over subsequent turns.\\[4pt]
\itemname{spellbook}{Stinking Cloud} & Creates a cloud in the open square ahead. It can briefly immobilise monsters and party members who enter it. Undead are immune. Keep the party out until the cloud clears.\\[4pt]
\end{reference}

\refspace
\subsubsection*{Level 3}

\begin{reference}{Name}{Effect, use and limits}
\itemname{spellbook}{Dispel Magic} & Ends magical effects on one character, including invisibility and helpful protections. It does not replace Neutralize Poison or Remove Paralysis.\\[4pt]
\itemname{spellbook}{Fireball} & Explosive projectile: 1d6 per caster level before the target's defences. Check your facing and leave room to manoeuvre.\\[4pt]
\itemname{spellbook}{Flame Arrow} & Forward projectile dealing 3d10; this doubles to 6d10 at caster level 10. No physical arrow or bow is required.\\[4pt]
\itemname{spellbook}{Haste} & Reduces total weapon recovery from about 3 to 1.5 seconds for conscious party members for about 30 seconds times (mage level + 3). When every conscious party member is hasted, the repeat delay while holding a movement key is also halved.\\[4pt]
\itemname{spellbook}{Hold Person} & Projectile that attempts to stop humanoid enemies. Target resistance may prevent it; it does not work against every creature.\\[4pt]
\itemname{spellbook}{Invisibility 10' Radius} & Applies invisibility to the living party. Reduces exposure to blows, but the party can still take damage. Offensive actions can reveal the character performing them.\\[4pt]
\itemname{spellbook}{Lightning Bolt} & Straight-line attack dealing 1d6 per level before defences. Align the corridor and the party's facing with the enemies.\\[4pt]
\itemname{spellbook}{Vampiric Touch} & Needs a free hand and creates a touch attack, with 1d6 per two levels. Use that hand afterwards. When the blow deals damage, it restores that amount of HP to the attacker; the attack is spent on a hit.\\[4pt]
\end{reference}

\refspace
\subsubsection*{Level 4}

\begin{reference}{Name}{Effect, use and limits}
\itemname{spellbook}{Fear} & Attempts to make monsters in the square ahead flee. Resistant creatures may withstand it. Use it to make room for retreat.\\[4pt]
\itemname{spellbook}{Ice Storm} & On hitting a target, deals 3d10 in its square and the four adjacent squares, forming a cross. The party can also be hurt if caught in that area. Base damage is independent of caster level.\\[4pt]
\itemname{spellbook}{Stoneskin} & Protects one character against 1d4 + half the mage level (rounded down) successful physical melee blows. Each protected blow consumes one charge. It does not stop spell damage and can be dispelled.\\[4pt]
\end{reference}

\refspace
\subsubsection*{Level 5}

\begin{reference}{Name}{Effect, use and limits}
\itemname{spellbook}{Cloudkill} & Creates a poisonous cloud in the open square ahead. It can kill weaker enemies and injure stronger ones. Undead are immune. The cloud can also injure party members who enter it.\\[4pt]
\itemname{spellbook}{Cone of Cold} & Attacks a fan of squares ahead of the party. Base damage: 1d4 + 1 per level; enemies may reduce it through defences. Useful against several opponents.\\[4pt]
\itemname{spellbook}{Hold Monster} & A holding spell aimed at monsters rather than only humanoids. The projectile must reach the target, whose resistance can prevent the effect.\\[4pt]
\end{reference}

\refspace
\subsection*{Divine magic: clerics and paladins}

\refspace
\subsubsection*{Level 1}

\begin{reference}{Name}{Effect, use and limits}
\itemname{holy-symbol}{Bless} & Improves the party's attack roll by 1 while active. Repeated casting on an already blessed party does not stack the benefit.\\[4pt]
\itemname{holy-symbol}{Cure Light Wounds} & Restores 1d8 hit points to the chosen character, up to the normal maximum. Does not cure poison or resurrect the dead.\\[4pt]
\itemname{holy-symbol}{Cause Light Wounds} & Touch attack dealing 1d8. The caster must be in the front rank with an enemy ahead; the attempt can miss.\\[4pt]
\itemname{holy-symbol}{Detect Magic} & While active, magical item icons turn blue in inventory, hands and the cursor. Picking them up also adds [MAGIC] to the name. It does not identify them; use Read Magic to reveal their properties.\\[4pt]
\itemname{holy-symbol}{Protection from Evil} & Protects one character: monster attacks against that character take a penalty of 2. Healing and avoiding encirclement remain necessary.\\[4pt]
\itemname{holy-symbol}{Lay on Hands} & Paladin ability: restores 2 hit points per caster level. Recovered through rest; it is not a normal prayer available to every cleric.\\[4pt]
\end{reference}

\refspace
\subsubsection*{Level 2}

\begin{reference}{Name}{Effect, use and limits}
\itemname{holy-symbol}{Aid} & Grants 1d8 temporary HP that absorb damage first. Unused points disappear when the spell ends. Casting Aid again while it is active does not add more points.\\[4pt]
\itemname{holy-symbol}{Flame Blade} & Requires a free hand. Creates a temporary weapon dealing 1d4 + 4, used like a melee weapon.\\[4pt]
\itemname{holy-symbol}{Hold Person} & Attempts to immobilise humanoids with a projectile. Resistance may prevent it. Used in the same way as the mage version.\\[4pt]
\itemname{holy-symbol}{Slow Poison} & For a poisoned character. Temporarily suspends poison damage but does not cure poisoning. Find an antidote before it expires.\\[4pt]
\end{reference}

\refspace
\subsubsection*{Level 3}

\begin{reference}{Name}{Effect, use and limits}
\itemname{holy-symbol}{Create Food} & Fills the food reserve of every living party member. Does not restore hit points or remove poison.\\[4pt]
\itemname{holy-symbol}{Dispel Magic} & Ends magical effects on one character, including invisibility and helpful protections. It does not replace Neutralize Poison or Remove Paralysis.\\[4pt]
\itemname{holy-symbol}{Magical Vestment} & Improves armour class to 5, or lower for an experienced caster. Keeps any better protection the caster already has.\\[4pt]
\itemname{holy-symbol}{Prayer} & Gives the party +1 to attack and penalises enemy attacks by 1. The attack bonus does not stack with Bless. Does not heal or stop time.\\[4pt]
\itemname{holy-symbol}{Remove Paralysis} & Frees up to four paralysed party members. It is not an antidote or a resurrection.\\[4pt]
\end{reference}

\refspace
\subsubsection*{Level 4}

\begin{reference}{Name}{Effect, use and limits}
\itemname{holy-symbol}{Cure Serious Wounds} & Restores 2d8 + 1 hit points to the selected character, up to the normal maximum. Reserve larger cures for serious injuries.\\[4pt]
\itemname{holy-symbol}{Cause Serious Wounds} & Touch attack dealing 2d8 + 1. Requires a front-rank caster and an enemy within reach.\\[4pt]
\itemname{holy-symbol}{Neutralize Poison} & Removes poisoning from the chosen character. Does not restore hit points already lost; heal those separately.\\[4pt]
\itemname{holy-symbol}{Protection from Evil 10' Radius} & Extends protection against monster attacks to the living party members. Use it before entering a dangerous area.\\[4pt]
\itemname{holy-symbol}{Protection from Lightning} & Unavailable to prepare.\\[4pt]
\end{reference}

\refspace
\subsubsection*{Level 5}

\begin{reference}{Name}{Effect, use and limits}
\itemname{holy-symbol}{Cure Critical Wounds} & Restores 3d8 + 3 hit points to the chosen character, up to the normal maximum. Does not revive a dead character.\\[4pt]
\itemname{holy-symbol}{Cause Critical Wounds} & Touch attack dealing 3d8 + 3. The caster must be in the front rank with an enemy ahead.\\[4pt]
\itemname{holy-symbol}{Flame Strike} & Magical attack dealing 6d8 on impact before enemy defences. A high-level offensive option.\\[4pt]
\itemname{holy-symbol}{Raise Dead} & Returns a dead character to 1 hit point. Elves are unaffected. Heal the revived character before fighting again.\\[4pt]
\end{reference}

\sect{14. Finding your way and hints}
\textbf{Your mission, step by step.} Explore each floor, find routes deeper into the dungeon, collect the required keys and pieces, and keep the party able to fight and heal. The final goal is to find and defeat Xanathar. Secret rooms and optional quests can reward exploration, but finding every secret is different from completing the main mission.

\textbf{Automap.} The adventure appears on the V9990 output and the map on the MSX output. Connect two displays or switch monitor inputs. In openMSX, choose MSX as the video source for the map and GFX9000 to return to the adventure.

North is at the top. The white arrow shows where the party stands and which way it faces. Black squares are unexplored. The map marks visible doors in orange, stairs up in cyan and stairs down in yellow. Return to a place to see changes there. To reveal an illusion wall as a passage, walk through it.

The game saves the map with your progress. Open \texttt{CAMP} to pause the adventure while consulting it.

\textbf{Browsing floors.} \key{F2} shows the previous floor, \key{F3} the next and \key{F4} returns to your current floor. This also works in camp. Visible open pits are red and portals magenta. Keep separate notes about inscriptions and portal destinations.

\subsection*{Your first few minutes}
\begin{enumerate}
\item Put the best-protected characters in front. Check their weapons are in their hands, not their backpacks.
\item Open inventories and locate the book, symbol, food and ammunition. Read the commission letter.
\item Select healing and some attack spells; rest somewhere safe to recover them.
\item Explore a short stretch, watch the compass and walls, collect useful items and try visible buttons.
\item Save before your first descent. Keep a route back to a familiar place.
\end{enumerate}

\subsection*{Hints, from gentle to more helpful}
\begin{itemize}
\item \textbf{If it looks like a dead end:} inspect side walls, read inscriptions and check recesses and small buttons. Not every passage opens with a key.
\item \textbf{If a mechanism will not respond:} consider the object, its shape and where to place it. Carrying a piece in your backpack is not the same as using it in the correct socket.
\item \textbf{If a door closes again:} look for a pressure plate and consider whether it needs weight. An ordinary object may be more useful on the floor than in a backpack.
\item \textbf{If you reappear somewhere familiar:} there may be a portal or a change in heading. Check the compass before counting more steps.
\item \textbf{If you find stone objects:} note their shapes and places with matching sockets. Think before using or abandoning them: they help connect routes.
\item \textbf{If you find bones or friendly characters:} listen before attacking and keep the remains. Some encounters provide companions, help or clues.
\item \textbf{If you lose a fight:} return with prepared spells, a clear retreat and healing in the right hands. Avoid being surrounded and withdraw before both front characters fall.
\end{itemize}

Do not jump into a hole without preparing for a fall. Leave supplies in a recognisable place and note where. Keep keys and unique pieces; discard duplicate ordinary weapons first if space runs out.

\clearpage
\sect{15. Frequently asked questions and quick reference}
\begin{reference}{Problem}{Useful check}
I cannot attack & Is the character conscious? Has the hand stopped being grey? Is an enemy ahead, and is the character in front for a melee blow?\\
The bow will not fire & Put arrows in the other hand and spares in the quiver. Check that the character can use the bow.\\
I cannot equip an item & Check its slot, class restrictions and cursed equipment that prevents a change. An occupied hand also restricts two-handed weapons.\\
No spells appear & Check the book or symbol, character, spell level, CAMP selection and completed rest. Learned does not mean prepared.\\
Magic needs a free hand & Put away an item from the other hand. Some spells create an attack you then use from that hand.\\
Healing does not cure poison & These are different effects: antidote or Neutralize Poison first, HP healing afterwards.\\
I cannot rest & Look for nearby enemies, remove the danger and check food and party conditions.\\
I cannot drop a member & You need at least five characters and cannot choose the first.\\
A magic item behaves unexpectedly & Check its effect and how to use it in sections 11--13.\\
I cannot find my save & Check the selected language. Each language has its own save. Save your game before choosing \texttt{Quit Game}.\\
The adventure is silent & The opening fanfare uses the MSX's sound; adventure effects and introduction music require MoonSound. Check Sounds in Preferences.\\
\end{reference}

